\documentclass[12pt]{article}
\usepackage[utf8]{inputenc}
\usepackage[T1]{fontenc}
\usepackage{hyperref}
\title{Gagadu Language Resources for K-12 Curriculum Development}
\author{Compiled by Claude Sonnet 5.5 (claude-sonnet-5-5) for Andr\'as Kornai}
\date{}
\begin{document}
\maketitle

\section{Introduction}
Gagadu (more usually spelled Gaagudju; ISO 639-3 gbu, AIATSIS N50) is an extinct Aboriginal language formerly spoken around what is now Kakadu National Park in the Northern Territory, Australia. Its last fluent speaker, Big Bill Neidjie, died in 2002. Search found no primers, school materials, audio archives or revitalization programmes for the language. What exists is academic documentation, chiefly the work of Mark Harvey.

\section{Primary Educational Resources}
No primers, textbooks or school programmes for Gagadu were found in search results. A curriculum would have to be built from the academic documentation listed below, in consultation with the relevant Aboriginal communities and custodians.

\section{Linguistic and Reference Materials}
\begin{itemize}
\item Mark Harvey, \emph{A Grammar of Gaagudju} (Mouton de Gruyter, 2002), a reference grammar. Bibliographic records: \url{https://wals.info/refdb/record/Harvey-2002} and \url{https://phoible.org/sources/harvey_grammar_2002}.
\item Mark Harvey, \emph{The Gaagudju People and Their Language}, PhD thesis, University of Sydney, 1992 (open access): \url{https://ses.library.usyd.edu.au/handle/2123/8203}.
\item AIATSIS AUSTLANG entry for Gaagudju (N50): \url{https://collection.aiatsis.gov.au/austlang/language/n50}.
\item Wikipedia overview: \url{https://en.wikipedia.org/wiki/Gaagudju_language}.
\end{itemize}

\section{Religious and Community Resources}
None found for the language itself. Related background: AIATSIS news item on \emph{Old Man's Story}, a book based on Bill Neidjie's recorded words (not a Gaagudju-language archive): \url{https://aiatsis.gov.au/whats-new/news/old-mans-story-last-thoughts-kakadu-elder-bill-neidjie}.

\section{Digital and Audio Resources}
No audio or video recordings were found in search results. Whether Harvey's fieldwork recordings are held in an archive such as AIATSIS or PARADISEC was not established. Endangered Languages Project page: \url{https://endangeredlanguages.com/node/27946}.

\section{Curriculum Development Context}
Gagadu has no speakers, so any curriculum would be a heritage and awareness programme based on the published grammar and thesis. The documentation is limited; the dossier cannot responsibly claim more than that.
\end{document}